L'éthanol pur est un liquide. Sa densité par rapport à l'eau $d=0,79$
et sa formule \ce{C2H5OH}.
\begin{enumerate}
    \item Calculer la quantité de matière de l'éthanol qui se trouve
          dans le volume $V=\qty{100}{\mL}$ de ce liquide.
    \item Déduire la masse de cette échantillon de l'éthanol.
\end{enumerate}

\begin{donnees}
    \begin{itemize}
        \item $M(O)=\qty{16}{\g\per\mol}$,
              $M(C)=\qty{12}{\g\per\mol}$, $M(H)=\qty{1}{\g\per\mol}$
        \item La masse volumique de l'eau~: $\rho_e=\qty{1}{\g\per\mL}$
    \end{itemize}
\end{donnees}
